\PuzzleChapterIntro{The Mirror Gate of Caldera}{The Obsidian Bisectors}{This is a coordinate geometry puzzle focusing on lines and angles. You must decode the inscription to find the primary lines, then calculate a sequence of angle bisectors based on specific geometric constraints.}
\Lore{
In Caldera’s basalt gate, three seams look as if they were drawn with light and then cooled into stone: two slanting rails and one vertical spine at \(x=1\).
Near the rails, someone has scratched the numbers \(7\) and \(13\) like a reminder left for the next worker.

A soot-stained slip has been wedged into the mortar at ankle height:

\begin{quote}
\textit{At the spine, take the sharp split. At the right side, take the wide. When a choice remains, keep the ray that points where it’s told.}
\end{quote}

Arin’s rule is simple: every “mirror” line demands an internal bisector, and the gate remembers the choices.
}

\Problem
In the usual $xy$--coordinate plane, the Gate's inscription is
\[
\bigl((x-y)^2-20(x-y)+91\bigr)(x-1)=0.
\]

\textbf{Direction-angle convention.}
The \emph{direction angle} of a \emph{ray} is the angle in $[0^\circ,360^\circ)$ measured counterclockwise from the positive $x$--axis to that ray.
The direction angle of a \emph{line} is defined to be the direction angle of its ray with positive $y$--component (so it lies in $[0^\circ,180^\circ)$; a vertical line has direction angle $90^\circ$).

The set of points satisfying the inscription is the union of three lines.
Call the vertical line the \textsf{Spine} $s$.
Call the two parallel slanted lines $\ell_{\odot}$ and $\ell_{\mathrm{M}}$ so that $\ell_{\odot}$ is closer to the origin.

Now build mirrors as follows (each mirror is a line through the indicated hinge-point):
\begin{itemize}[leftmargin=1.2em]
\item At $\ell_{\odot}\cap s$, place mirror $h$ along the \emph{internal} bisector of the \emph{acute} angle between $\ell_{\odot}$ and $s$.
\item At $\ell_{\mathrm{M}}\cap s$, look in the half-plane $x>1$ and take the \emph{obtuse} angle formed there by the rays of $\ell_{\mathrm{M}}$ and $s$; place mirror $g$ along its \emph{internal} bisector.
\item Mirror $h$ meets $\ell_{\mathrm{M}}$ at a point away from the Gate. There, place mirror $j$ along the \emph{internal} angle bisector of $\ell_{\mathrm{M}}$ and $h$ \emph{whose ray with positive $y$--component points into the half-plane $x>1$}.
\item Mirrors $j$ and $g$ intersect; place the final mirror $m$ along the \emph{internal} angle bisector of $j$ and $g$ \emph{whose ray with positive $y$--component points into the half-plane $x<1$}.
\end{itemize}

Let $\delta$ be the direction angle of mirror $m$ (in the line sense above).

\VaultDirective{Determine $\delta$ as a reduced rational multiple of $\pi$.}

\SolutionPage
\small
Factor the inscription using \(t=x-y\):
\[
(t^2-20t+91)(x-1)=0=(t-7)(t-13)(x-1),
\]
so \(\ell_{\odot}:y=x-7\), \(\ell_{\mathrm{M}}:y=x-13\) (both direction \(45^\circ\)), and \(s:x=1\) (direction \(90^\circ\)).

\smallskip
\(h\): acute bisector of \(45^\circ\) and \(90^\circ\), so \(\theta_h=67.5^\circ=\frac{135}{2}^\circ\).

\(g\): in \(x>1\), the obtuse angle is between the \(45^\circ\) ray of \(\ell_{\mathrm{M}}\) and the downward vertical ray \(270^\circ\), of size \(135^\circ\).
Its bisecting ray has direction \(270^\circ+67.5^\circ=337.5^\circ\); as a \emph{line} we take the opposite direction with positive \(y\)-component:
\(\theta_g=337.5^\circ-180^\circ=157.5^\circ=\frac{315}{2}^\circ\).

\(j\): internal bisectors between \(45^\circ\) and \(67.5^\circ\) are \(56.25^\circ\) and \(146.25^\circ\); the \(x>1\) rule selects
\(\theta_j=56.25^\circ=\frac{225}{4}^\circ\).

\(m\): internal bisectors between \(56.25^\circ\) and \(157.5^\circ\) are
\[
\frac{56.25^\circ+157.5^\circ}{2}=106.875^\circ \quad\text{and}\quad 16.875^\circ;
\]
the \(x<1\) rule selects \(\delta=106.875^\circ=\frac{855}{8}^\circ\).
Thus
\[
\delta=\frac{855}{8}\cdot\frac{\pi}{180}=\frac{19\pi}{32}.
\]

\begin{center}
\includegraphics[width=0.55\linewidth]{figures/p02_1.png}
\end{center}

\FinalAnswer{%
\(\delta=\dfrac{19\pi}{32}\).
}

\NextPage
